\chapter{Design and Methodology}\label{ch:B}

\section{System Architecture Overview}

The system consists of three main components:

\begin{enumerate}
  \item \textbf{Data pipeline} — dataset curation, physics-based augmentation,
        feature engineering, and model training (\texttt{train.py},
        \texttt{expand\_dataset.py}, \texttt{utils/preprocessing.py}).
  \item \textbf{Backend API} — a Flask REST service that loads trained model
        artifacts and serves predictions with confidence scores
        (\texttt{app.py}).
  \item \textbf{Frontend} — a single-page HTML/CSS/JavaScript application
        that collects material parameters and displays predicted results
        interactively.
\end{enumerate}

\section{Dataset}

\subsection{Collection}

The raw dataset was manually curated from 529 publications retrieved through
Web of Science and Google Scholar using the query terms
\emph{``photoelectrochemical water splitting''},
\emph{``photoanode photocurrent density''}, and
\emph{``STH efficiency semiconductor''}.
Each row corresponds to a single experimental condition reported in one paper,
identified by DOI.

The 30 original columns span material identity, structural characteristics,
electrochemical conditions, and performance outputs as listed in
Table~\ref{tab:columns}.

\begin{table}[h]
\centering
\caption{Dataset columns and coverage in the original (529-row) dataset.}
\label{tab:columns}
\begin{tabular}{lll}
\toprule
Column & Type & Coverage (\%) \\
\midrule
Material System         & Categorical & 97 \\
Bandgap (eV)            & Numeric     & 49 \\
Synthesis Method        & Categorical & 80 \\
Nanostructure           & Categorical & 67 \\
Cocatalyst              & Categorical & 62 \\
Electrolyte             & Categorical & 73 \\
pH                      & Numeric     & 47 \\
Light Intensity (mW/cm\textsuperscript{2}) & Numeric & 61 \\
Applied Bias (V vs RHE) & Numeric     & 74 \\
Temperature (K)         & Numeric     & 21 \\
Photocurrent Density    & Numeric     & 78 \\
STH Efficiency (\%)     & Numeric     & \textbf{9} \\
H\textsubscript{2} Evolution Rate & Numeric & \textbf{11} \\
\bottomrule
\end{tabular}
\end{table}

The extreme sparsity of STH (9\%) and H\textsubscript{2} (11\%) is the
central data challenge: with only 48 and 57 usable rows respectively,
5-fold cross-validation operates on as few as 10 validation samples per fold,
making R\textsuperscript{2} estimates highly variable.

\subsection{Physics-Based Dataset Augmentation}

To address the sparsity problem, two complementary augmentation strategies
were applied, producing a final dataset of \textbf{676 rows}.

\subsubsection{Strategy 1: Physics-Derived Estimates for Existing Rows}

For the 364 rows that have a measured photocurrent but no reported STH,
the STH value was estimated using the standard approximation:

\begin{equation}
  \widehat{\text{STH}} = \frac{J_{ph} \times 1.23}{P_{in}} \times \eta_F \times 100\%
  \label{eq:sth_est}
\end{equation}

where $P_{in}$ is set to the reported light intensity when available,
or to 100\,mW/cm\textsuperscript{2} (1 sun, AM~1.5G standard) otherwise;
$\eta_F$ is sampled uniformly from $[0.75,\,0.98]$ (unbiased measurement) or
$[0.55,\,0.85]$ (biased measurement, $|V_{bias}| > 0.1$\,V) to reflect realistic
Faradaic losses.
A multiplicative noise factor $\delta \sim \mathcal{U}(0.92,\,1.08)$ is added
to reproduce typical measurement variability.

Similarly, H\textsubscript{2} evolution rate was derived from Faraday's law
(Equation~\ref{eq:her}) with $\eta_F \sim \mathcal{U}(0.60,\,0.95)$.

This strategy filled \textbf{361 STH} and \textbf{357 H\textsubscript{2}} values.
The approach is scientifically justified because the relationship
$\text{STH} \propto J_{ph}$ is a thermodynamic identity, not a model assumption.
The introduced noise prevents the augmented labels from being perfect linear
functions of the input features, preserving the learning signal.

\subsubsection{Strategy 2: Literature-Based Synthetic Rows}

An additional 147 rows were generated for 25 well-characterised photoelectrode
systems (BiVO\textsubscript{4} variants, $\alpha$-Fe\textsubscript{2}O\textsubscript{3},
WO\textsubscript{3}, TiO\textsubscript{2}, ZnO, Cu\textsubscript{2}O, GaN,
Ta\textsubscript{3}N\textsubscript{5}, Si, g-C\textsubscript{3}N\textsubscript{4},
CuInS\textsubscript{2}, and others).
Photocurrent ranges for each material--cocatalyst combination were taken from
published review articles \cite{dias2016recent,tamirat2016using}.
All structural parameters (electrolyte, pH, illumination) match standard
experimental conditions for the corresponding material.
Every synthetic row is tagged \texttt{Data\_Source = "synthetic"} in the dataset.

\subsection{Missing Value Imputation}

Before augmentation, additional standard values were filled deterministically:
\begin{itemize}
  \item \textbf{Bandgap}: for 204 rows with known material but missing bandgap,
        reference values were filled from a lookup table of 40+ materials
        (e.g., TiO\textsubscript{2} anatase: 3.2\,eV, BiVO\textsubscript{4}: 2.4\,eV,
        Fe\textsubscript{2}O\textsubscript{3}: 2.1\,eV) with a small jitter
        $\delta \sim \mathcal{U}(-0.05,\,+0.05)$\,eV.
  \item \textbf{pH}: inferred from electrolyte composition for 114 rows
        (KOH/NaOH $\to$ pH 13--14, H\textsubscript{2}SO\textsubscript{4} $\to$ pH 0--1.5,
        Na\textsubscript{2}SO\textsubscript{4} $\to$ pH 5.8--7.2, etc.).
  \item \textbf{Temperature}: set to 298\,K for all 420 rows lacking this value
        (standard laboratory conditions).
  \item \textbf{Light intensity}: set to 100\,mW/cm\textsuperscript{2} for 61 rows
        with AM~1.5G illumination but missing numeric intensity.
\end{itemize}

\subsection{Final Dataset Statistics}

After augmentation, the dataset contains 676 rows with the coverage shown in
Table~\ref{tab:augmented}.

\begin{table}[h]
\centering
\caption{Dataset coverage before and after augmentation.}
\label{tab:augmented}
\begin{tabular}{lrrrr}
\toprule
Target / Feature & Before & Before (\%) & After & After (\%) \\
\midrule
Photocurrent Density  & 410 & 78 & 557 & 82 \\
STH Efficiency        & 48  & 9  & 556 & 82 \\
H\textsubscript{2} Evolution Rate & 57 & 11 & 559 & 83 \\
Bandgap               & 258 & 49 & 609 & 90 \\
pH                    & 250 & 47 & 509 & 75 \\
Temperature           & 112 & 21 & 676 & 100 \\
\bottomrule
\end{tabular}
\end{table}

\section{Feature Engineering}

Twenty numeric features are derived from eight raw input fields.
Table~\ref{tab:features} lists all engineered features with their physical
motivation.

\begin{table}[h]
\centering
\caption{Physics-informed engineered features.}
\label{tab:features}
\begin{tabular}{lp{7cm}}
\toprule
Feature & Physical motivation \\
\midrule
\texttt{photon\_energy} & $1240/E_g$ — absorption edge wavelength (nm) \\
\texttt{overpotential\_proxy} & $V_{bias} - (E_g - 1.23)$ — net electrochemical driving force beyond thermodynamic minimum \\
\texttt{bandgap\_sq} & $E_g^2$ — non-linear dependence of absorption on gap \\
\texttt{log\_bandgap} & $\ln(1+E_g)$ — compressed bandgap scale \\
\texttt{ph\_deviation} & $|pH - 7|$ — distance from neutral (surface charge proxy) \\
\texttt{bias\_positive} & $\max(V_{bias},\,0)$ — forward bias contribution only \\
\texttt{temp\_normalized} & $(T - 298)/298$ — deviation from room temperature \\
\texttt{has\_cocatalyst} & Binary: cocatalyst present \\
\texttt{is\_nanostructured} & Binary: nanostructure vs.\ bulk thin film \\
\texttt{ionic\_strength\_proxy} & Mapped from electrolyte name (KOH=1.0, Na\textsubscript{2}SO\textsubscript{4}=0.6, PBS=0.5) \\
\texttt{bandgap\_x\_ph} & $E_g \times \text{pH}$ — flat-band potential shifts $-59$\,mV/pH (Nernst) \\
\texttt{bias\_x\_ionic} & $V_{bias} \times I_{proxy}$ — electric field in double layer \\
\texttt{bg\_x\_overpot} & $E_g \times \eta_{proxy}$ — combined material+electrochemical driving force \\
\texttt{ph\_x\_temp} & $\text{pH} \times T$ — Arrhenius-like surface kinetics \\
\bottomrule
\end{tabular}
\end{table}

Six categorical features (material system, electrolyte, synthesis method,
photoelectrode type, nanostructure, morphology) are one-hot encoded inside
a scikit-learn \texttt{Pipeline} with \texttt{handle\_unknown="ignore"}
so that unseen material names at inference time do not cause failures.

\section{Data Leakage Prevention}

STH is defined by Equation~\ref{eq:sth_est} as a deterministic function of
photocurrent.
Including photocurrent as a predictor for STH would inflate R\textsuperscript{2}
to near 1.0 without the model learning any generalisable relationship.
The STH model therefore excludes photocurrent from its feature set
(\texttt{TARGET\_EXCLUSIONS["sth"] = ["photocurrent"]}).
A Pearson correlation audit is logged at training time to confirm this leakage is
absent.

Similarly, 33 photocathode entries (negative photocurrent by measurement convention)
are removed from the photocurrent training set to avoid a bimodal target
distribution that would confuse regression models.

\section{ML Model Pipeline}

\subsection{Model Selection}

Two tree ensemble methods are trained per target:

\begin{enumerate}
  \item \textbf{XGBoost Regressor} — regularised gradient boosting.
        Hyperparameters: 300 trees, learning rate 0.05, max depth 4,
        subsample 0.8, colsample\_bytree 0.8, $L_1 = 0.1$, $L_2 = 1.0$,
        min child weight 3.

  \item \textbf{Random Forest Regressor} — bagged decision trees.
        Hyperparameters: 300 trees, max depth 8, min samples leaf 2,
        max features ``sqrt''.
\end{enumerate}

Both models are wrapped in a scikit-learn \texttt{Pipeline} with a
\texttt{ColumnTransformer} preprocessor (median imputation for numerics,
mode imputation + OHE for categoricals).
The full pipeline is cloned and re-fit on all training data after CV evaluation
to produce the inference artifact.

\subsection{Ensemble Inference and Confidence Scoring}

At prediction time, both saved models are queried.
The ensemble value is the CV-R\textsuperscript{2}-weighted average:

\begin{equation}
  \hat{y}_{ens} = \frac{\sum_m \max(0,\,r^2_m)\,\hat{y}_m}{\sum_m \max(0,\,r^2_m)}
\end{equation}

Prediction uncertainty is characterised by the \emph{relative standard deviation}
of the two models:
\begin{equation}
  \sigma_{rel} = \frac{\text{std}(\hat{y}_1,\hat{y}_2)}{\hat{y}_{ens}}
\end{equation}

The \emph{confidence score} combines model quality and agreement:
\begin{equation}
  c = 0.6 \times \overline{r^2} + 0.4 \times \max(0,\,1 - \sigma_{rel})
  \label{eq:confidence}
\end{equation}
where $\overline{r^2}$ is the mean CV R\textsuperscript{2} clipped to $[0,1]$.
A score of 1.0 would require both a perfect model and zero disagreement between
the two estimators.

\subsection{Performance Classification}

The composite performance score is:
\begin{equation}
  s = \left(p_{pc} + b_{STH}\right) \times (0.4 + 0.6\,c)
  \label{eq:score}
\end{equation}
where $p_{pc} \in [0,\,100]$ is the photocurrent percentile position relative to
the training distribution (p75 = 100 points, linear interpolation between p40
and p75), $b_{STH} \in \{0, 7, 15\}$ is a STH bonus, and $c$ is the confidence
from Equation~\ref{eq:confidence}.
The confidence factor $(0.4 + 0.6\,c)$ downgrades speculative predictions:
a model with $r^2 = 0$ contributes at most 40 points.
Final labels are HIGH ($s \geq 70$), MEDIUM ($40 \leq s < 70$), LOW ($s < 40$).

\subsection{Cross-Validation Protocol}

The number of CV folds is capped at $\min(5,\,\lfloor n/5 \rfloor)$ to prevent
degenerate folds on sparse targets.
For the augmented dataset (556 STH rows), full 5-fold CV is used.
The final model is fit on all available data for each target.

\section{System Implementation}

\subsection{Training Pipeline}

The training script (\texttt{train.py}) follows four stages:
\begin{enumerate}
  \item Load and map raw CSV column names to internal identifiers.
  \item Apply physics-based cleaning (numeric parsing, sanity bounds).
  \item Engineer all 20 features via \texttt{utils/preprocessing.py}.
  \item Train both models per target, evaluate with CV, save the artifact
        as a \texttt{joblib} pickle containing both pipelines,
        CV metrics, feature lists, training medians/modes, and
        percentile thresholds.
\end{enumerate}

\subsection{REST API}

The Flask application exposes three endpoints:
\begin{itemize}
  \item \texttt{GET /} — health check; returns loaded model names and CV R\textsuperscript{2}.
  \item \texttt{POST /predict} — accepts a JSON body with material parameters,
        returns predictions for all three targets with confidence and
        performance classification.
  \item \texttt{GET /dataset} — returns the clean training dataset as JSON
        for the frontend data explorer.
\end{itemize}

Input validation enforces physical bounds (pH $\in [0,14]$, bandgap $\in [0.5,6]$\,eV,
bias $\in [-2,5]$\,V vs.\ RHE) and raises 400 errors with descriptive messages.
All errors are logged with stack traces via Python's \texttt{logging} module.

\subsection{Web Application}

The frontend is a single HTML file served directly by Flask.
It provides:
\begin{itemize}
  \item A \textbf{Quick Predict} tab with a minimal 4-field form (material,
        electrolyte, pH, bias) for rapid prototyping.
  \item A \textbf{Full Predict} tab with all 15 optional parameters.
  \item A \textbf{Results} tab displaying values, model confidence, a
        colour-coded performance bar, and input summary.
  \item A \textbf{Dataset} explorer showing the training data with
        sortable columns.
\end{itemize}
